\documentclass[11pt,reqno]{amsart}

%==============================================================
% Does Random-Clock, Heavy-Tailed Structure Compress?
% An empirical companion to the iterated-Brownian / operator papers.
%==============================================================

\usepackage[utf8]{inputenc}
\usepackage[T1]{fontenc}
\usepackage{lmodern}
\usepackage[margin=1.15in]{geometry}
\usepackage{amsmath,amssymb,amsthm,mathtools}
\usepackage{microtype}
\usepackage{booktabs}
\usepackage{hyperref}
\usepackage{cleveref}

\theoremstyle{plain}
\newtheorem{theorem}{Theorem}[section]
\newtheorem{proposition}[theorem]{Proposition}
\theoremstyle{definition}
\newtheorem{definition}[theorem]{Definition}
\theoremstyle{remark}
\newtheorem{remark}[theorem]{Remark}
\newtheorem*{finding}{Finding}

\newcommand{\R}{\mathbb{R}}
\newcommand{\E}{\mathbb{E}}
\newcommand{\IBM}{Z}
\newcommand{\bps}{\mathrm{bits/sample}}
\DeclareMathOperator{\Var}{Var}

\title[Does random-clock, heavy-tailed structure compress?]{Does Random-Clock, Heavy-Tailed Structure Compress?\\
  An Empirical Study of Iterated-Brownian and\\ Operator Models for Lossless Coding}
\author{Ramiro Fontes}
\address{Quijotic Research}
\email{ramirofontes@gmail.com}
\urladdr{www.quijotic.com}
\date{\today}

\begin{document}

\begin{abstract}
The companion papers develop, for iterated Brownian motion (IBM)
$\IBM_t=X_{|Y_t|}$ and the order-$p$ heat pseudo-processes, a heavy-tailed
marginal calculus (operator order $p$, biharmonic tail
$e^{-c|r|^{4/3}}$), a long-memory volatility structure (fractional /
Mittag-Leffler), and a shared adjoint operator framework $D=\delta^{*}$.
This note asks the only question that decides whether any of that helps
\emph{lossless compression}: can these models lower the held-out code length
$-\tfrac1n\sum_i\log_2 q_i(x_i\mid x_{<i})$ below standard baselines? We run a
five-stage study---synthetic go/no-go, real-data tail coding, a parameter-swept
Mittag-Leffler volatility model against EWMA/GARCH, a byte-level arithmetic-coded
codec against gzip/bzip2/xz and a FLAC-class baseline, and a GPU neural
distribution-head model---on synthetic IBM, a real electrocardiogram
($108\text{k}$ samples), pooled S\&P~500 daily log-returns ($619\text{k}$), and
$16$-bit speech. The result is a clear but qualified yes. Non-Gaussian
residual coding gives a real, measurable advantage: $+0.07$ to $+0.23~\bps$
over Gaussian on real data, beating the Laplace/Rice coder of FLAC. Tested as a
real codec---canonical raw bytes, exact byte-for-byte round-trip, all overhead
included---it beats gzip, bzip2, and xz on numeric streams ($+8\%$ to $+31\%$
over xz) and \emph{ties} a faithful FLAC-class (LPC$+$Rice) baseline on
continuous $16$-bit speech ($+0.5\%$, within real-FLAC tuning margin). The
advantage is largest under a learned predictor, where a heavy-tailed
distribution head is worth $+0.17$--$0.20~\bps$ over a Gaussian head at equal
model size. The sobering byte-level lesson is that the \emph{driver is
prediction plus arithmetic/Rice coding}: a plain Gaussian model captures most of
the margin, the heavy-tail and clock layers add only $\approx1$--$3\%$ on top,
and on audio every static heavy-tail likelihood (the IBM mixture included) loses
to Rice coding---only the volatility model is a consistent lever. Two predictions of the theory are
\emph{falsified} on real data, however: the winning tail is power-law
(Student-$t$), not the operator's stretched-exponential $\beta=4/3$; and the
fractional volatility kernel, while confirming genuine long memory
($\alpha<1$ optimal on every dataset), beats GARCH only marginally and not
universally. We test the genuine objects directly: the true IBM mixture (a
half-normal variance mixture, not a tail proxy) is the best of the four
canonical codes on synthetic IBM residuals but shares the stretched-exponential
tail that loses to Student-$t$ on real data; and the raw order-$p$
pseudo-process kernels are \emph{signed} ($12\%$ negative mass for the
biharmonic, wildly oscillatory for Airy), hence not entropy codes at all---only
their positive subordinated (IBM) face compresses. The durable contribution for
a codec is therefore the
\emph{framework}---adaptive heavy-tail selection and a jointly modeled
volatility clock---rather than the specific operator exponents.
\end{abstract}

\maketitle

\section{The compression question}
\label{sec:intro}

A lossless compressor is a predictive model followed by an entropy coder. The
coder (arithmetic, range, or ANS) is solved: it realizes the model's code
length to within a fraction of a bit. Hence the entire question is whether a
model's predictive distribution $q_i(x_i\mid x_{<i})$ assigns shorter ideal
code length
\begin{equation}\label{eq:codelen}
  L = -\frac1n\sum_{i=1}^{n}\log_2 q_i(x_i\mid x_{<i})\quad(\bps)
\end{equation}
than a baseline. For continuous residuals \eqref{eq:codelen} is a differential
code length; every model pays the same quantization-bin constant, which cancels
in model-versus-model differences, so relative $\bps$ is the honest comparison.

The companion papers suggest three mechanisms that might lower
\eqref{eq:codelen} on bursty signals.
\begin{enumerate}
\item \emph{Heavy tails.} The IBM marginal is a normal variance mixture; by a
  Laplace-method estimate its tail is the biharmonic
  $q(z)\sim\exp(-c|z|^{4/3})$, i.e.\ exponential-power with shape
  $\beta=4/3$. The order-$p$ pseudo-process generalizes this to
  $\beta=p/(p-1)$ (Gaussian $p{=}2\!\to\!\beta{=}2$; biharmonic
  $p{=}4\!\to\!\beta{=}4/3$).
\item \emph{A volatility clock.} The random time-change makes the local
  variance vary; if predictable, it shortens the code (volatility clustering).
\item \emph{Long memory.} The fractional / Mittag-Leffler structure of the
  memory companion predicts power-law, not exponential, volatility
  autocorrelation.
\end{enumerate}
The order-$p$ Stein identity (Stein operator $=$ kernel score) further suggests
a closed-form, fit-free selector for the tail exponent.

We test these as falsifiable claims. Throughout, the headline metric is
held-out $\bps$; the kill criterion is whether the structure beats the
appropriate baseline by a margin that survives on real data.

\section{Method and data}
\label{sec:method}

\paragraph{Datasets.} (i) \emph{Synthetic IBM}: genuine iterated Brownian
increments from a fixed Brownian scenery sampled at a reflected Brownian clock
($2\times10^5$ samples), which carries both the heavy tail and the true
long-range dependence. (ii) \emph{ECG}: the SciPy electrocardiogram, a real
$108{,}000$-sample biosignal with bursty QRS complexes. (iii) \emph{Equities}:
daily log-returns pooled across $505$ S\&P~500 names, $618{,}535$ samples.

\paragraph{Models.} Residuals are formed by an order-$p$ AR predictor (rounded
to integers where lossless coding is required). The residual code is one of:
Gaussian, Laplace ($\approx$ the FLAC/Rice coder), Student-$t$, the
\emph{genuine IBM mixture} (a normal variance mixture with half-normal mixing,
the true IBM-increment marginal, one scale parameter), exponential-power with
fitted or operator-fixed shape $\beta=p/(p-1)$, and dynamic variants with an
EWMA, GARCH$(1,1)$, power-law, or genuine Mittag-Leffler volatility scale. Parameters are fit on a training
prefix; $\bps$ is reported on a disjoint held-out span.

\paragraph{Diagnostics.} The fitted tail exponent $\hat\beta$ and its implied
operator order $p=\beta/(\beta-1)$; the lag-1 autocorrelation of squared
residuals $\mathrm{ACF}_1(r^2)$ (presence of clustering); and the decay law of
$\mathrm{ACF}_k(r^2)$ (power-law $=$ long memory vs.\ exponential $=$ short).

\section{Stage 0: synthetic go/no-go}
\label{sec:stage0}

On genuine IBM the best static heavy-tail code beats Gaussian by
$+0.049~\bps$, and the \emph{parameter-free} $\beta=4/3$ code ($2.025~\bps$)
beats fitted Student-$t$ ($2.032$). Two independent estimators of the tail
exponent agree, $\hat\beta=1.289$ (MLE) $=1.289$ (score-matching), next to the
predicted $4/3$. The Gaussian control returns $\hat\beta=1.99$ with no spurious
gain, so the machinery does not hallucinate tails.

Coding the \emph{genuine} IBM mixture density---a normal variance mixture with
half-normal mixing, one scale parameter, not the exponential-power proxy---against
the three canonical codes makes the load-bearing comparison
(\cref{tab:fourway}). On data that is actually IBM the mixture is the best of the
four, beating Student-$t$ by $\approx0.01~\bps$ on residuals it itself generated,
and it correctly \emph{loses} on Gaussian and power-law (stochastic-volatility)
data---it is the right code specifically for IBM-structured residuals, not a
universal hammer.

\begin{table}[h]
\centering\small
\begin{tabular}{lcc}
\toprule
data & winner & IBM mix vs.\ Gauss / Student-$t$\\
\midrule
true IBM-increment law & \textbf{IBM mixture} & $+0.058$ / $+0.009$\\
genuine IBM            & \textbf{IBM mixture} & $+0.051$ / $+0.011$\\
Gaussian control       & Gaussian             & $-0.040$ (correctly loses)\\
power-law SV           & Student-$t$          & (correctly loses)\\
\bottomrule
\end{tabular}
\caption{The genuine IBM mixture against the canonical codes (held-out $\bps$
gain). It wins on IBM-type residuals and correctly loses elsewhere.}
\label{tab:fourway}
\end{table}

The squared-residual autocorrelation of IBM decays as a power law,
$\mathrm{ACF}_k(r^2)\sim k^{-1/2}$ ($R^2_{\text{pow}}=0.87$ vs.\
$R^2_{\text{exp}}=0.69$), i.e.\ \emph{long memory}; an exponential
stochastic-volatility control decays exponentially, as expected. Consistently,
a power-law volatility filter recovers more of the clustering than EWMA
($+0.028$ vs.\ $+0.014~\bps$) on IBM, while EWMA wins on the exponential
control---the volatility tool must match the memory law.

\begin{finding}
The operator's specific tail prediction ($\beta\approx4/3$) holds on data
generated by the operator; clustering is present and is long-memory; both are
measurable in $\bps$.
\end{finding}

A disentanglement experiment (coding IBM, a non-IBM fractional process, and
i.i.d.\ heavy-tail data with one generic ``fractional$+$tail'' model against a
clock-aware oracle) showed that IBM, Mittag-Leffler, and the heavy tail are not
competing explanations but \emph{the same subordination structure seen twice}:
one random clock produces both the variance-mixture tail and the fractional
memory (the Allouba--Zheng subordination/PDE correspondence). The oracle gap on
IBM is large ($+0.21~\bps$), but it is mostly \emph{unpredictable-from-the-past}
variance, so the causally bankable clock gain is modest.

\section{The pseudo-process is signed; IBM is its codeable face}
\label{sec:signed}

We built the genuine order-$p$ kernels
$K_p(t,x)=\mathrm{IFFT}\bigl[e^{c_p(ik)^p t}\bigr]$ directly from the
characteristic symbol (paper~3, Rmk.~5.7); the reconstruction is exact---$K_2$
reproduces the standard Gaussian to $\|K_2-\mathcal N(0,1)\|_\infty=1.7\times
10^{-16}$. The central fact for compression is that an entropy code needs a
\emph{nonnegative} density: $-\log_2 q$ is undefined where $q<0$. The order-$p$
kernels fail this for $p\ge3$ (\cref{tab:signed}).

\begin{table}[h]
\centering\small
\begin{tabular}{lcccl}
\toprule
$p$ & $\int K$ & $\|K\|_{L^1}$ & neg.\ mass & status\\
\midrule
$2$ (Gaussian)   & $1.000$ & $1.000$  & $0$     & positive --- codeable\\
$4$ (biharmonic) & $1.000$ & $1.237$  & $0.119$ & \textbf{signed --- not a code}\\
$3$ (Airy)       & $1.000$ & $18.08$  & $8.54$  & \textbf{signed, near-non-integrable}\\
\bottomrule
\end{tabular}
\caption{Genuine order-$p$ kernels by FFT of the symbol. The biharmonic puts
$12\%$ of its mass in negative lobes; the Airy kernel is wildly oscillatory
($\|K\|_{L^1}=18$). For $p\ge3$, $K_p$ is signed and cannot be an entropy code.}
\label{tab:signed}
\end{table}

On genuine IBM residuals, $0.59\%$ fall where $K_4\le0$: the raw biharmonic code
length is literally undefined there. The best positive salvage, $|K_4|/\|K_4\|_{L^1}$
(a proper density), gives $2.139~\bps$---worse than every baseline---because
taking absolute values destroys the structure. By contrast the positive
subordinated law of the same order-$4$ operator, the \emph{IBM mixture}, gives
$2.019~\bps$ and beats Student-$t$ ($2.029$). ``Operator order $4$'' and ``IBM
marginal'' are the signed and positive sides of one object: the pseudo-process is
the signed PDE Green's function and does not compress; its IBM subordination is a
genuine probability and does.

\section{Stage 1: does the tail survive on real data?}
\label{sec:stage1}

\begin{table}[h]
\centering\small
\begin{tabular}{lcc}
\toprule
quantity & ECG (AR resid) & S\&P returns \\
\midrule
best static vs.\ Gaussian            & $+0.068$ & $+0.226$ \\
winning family                        & Student-$t$ & Student-$t$ \\
exp-power $\beta=4/3$ (operator)      & $+0.043$ & $+0.177$ \\
vs.\ Laplace (FLAC/Rice)             & $+0.025$ & $+0.015$ \\
with EWMA clustering                  & $+0.111$ & $+0.264$ \\
$\hat\beta$ (MLE)                     & $1.18$   & $1.05$ \\
$\mathrm{ACF}_1(r^2)$                 & $0.227$  & $0.091$ \\
\bottomrule
\end{tabular}
\caption{Held-out $\bps$ improvements over Gaussian on real data. The heavy
tail and clustering both pay; Student-$t$ wins; the operator $\beta=4/3$ is
competitive but not the winner.}
\label{tab:stage1}
\end{table}

The heavy-tail advantage is real and, on equities, large: $+0.226~\bps$ is
about $11\%$ below Gaussian residual coding, and it beats the Laplace coder
that FLAC uses. But \emph{Student-$t$ wins on both datasets}: real residuals are
power-law tailed (fitted $\mathrm{df}\approx3$), which is heavier than any
stretched exponential. The operator predicts the lighter $e^{-c|r|^{4/3}}$ and
therefore undershoots the true tail.

\begin{finding}
The generic claim ``non-Gaussian heavy tails compress better'' is strongly
confirmed on real data; the specific operator shape $\beta=4/3$ is falsified
there---real tails are power-law.
\end{finding}

\section{Stage 2: long-memory volatility (Mittag-Leffler) vs.\ GARCH}
\label{sec:stage2}

We tested the genuine Mittag-Leffler kernel
$w_k=k^{\alpha-1}E_{\alpha,\alpha}(-\lambda k^{\alpha})$ (the fractional-OU
volatility memory, with $\alpha=2H_2$), swept over $\alpha\in\{0.4,0.6,0.8,1.0\}$
and five $\lambda$, against EWMA, a fitted GARCH$(1,1)$, and a power-law filter,
all feeding a common Student-$t$ tail. The Mittag-Leffler function was computed
in high precision (series and asymptotic), validated to $10^{-7}$ against
$E_1(-x)=e^{-x}$.

\begin{table}[h]
\centering\small
\begin{tabular}{lccccc}
\toprule
dataset & Mittag-Leffler & GARCH$(1,1)$ & EWMA & static & optimal $\alpha$\\
\midrule
synthetic IBM      & $+0.032$ & $+0.027$ & $+0.018$ & $0$ & $0.4$\\
S\&P returns       & $+0.067$ & $+0.062$ & $+0.055$ & $0$ & $0.4$\\
ECG residual       & $+0.090$ & $+0.110$ & $+0.049$ & $0$ & $0.8$\\
\bottomrule
\end{tabular}
\caption{Volatility-model $\bps$ gain over a static scale (common Student-$t$
tail). The optimal fractional order is $\alpha<1$ on every dataset.}
\label{tab:stage2}
\end{table}

The optimum is $\alpha<1$ everywhere: volatility genuinely has long memory,
including in real financial and ECG data---a clean, literature-consistent
confirmation of the memory companion. Mittag-Leffler beats EWMA on all three
datasets and beats GARCH on the financial and IBM data; GARCH wins on ECG,
whose volatility is more persistent/short-memory. The Mittag-Leffler kernel and
a plain power-law kernel tie, since $E_{\alpha,\alpha}(-x)$ with $\alpha<1$ is
itself power-law tailed.

\begin{finding}
Long memory in volatility is real and the fractional kernel captures it, but its
edge over GARCH is small ($\sim0.005~\bps$ on equities) and not universal.
\end{finding}

\section{Stage 3: a real byte-level codec vs.\ gzip / bzip2 / xz}
\label{sec:stage3}

This stage is tested as a compressor, not a paper. The input is frozen as
canonical little-endian raw bytes (int16 for ECG, int32 scaled returns for
S\&P); every method sees exactly those bytes. Our codec must
$\textsf{encode}\to\textsf{decode}\to$ reproduce the input \emph{byte for byte}
(verified), and its compressed file carries \emph{all} overhead---magic, length,
dtype, scale, AR coefficients, model parameters---so it is fully self-decoding.
The coder is a Witten--Neal--Cleary arithmetic coder over a per-symbol
cumulative-frequency table built from the model CDF with a guaranteed nonzero
floor. Baselines are actual bytes from gzip\,$-9$, bzip2\,$-9$, xz\,$-9$
(LZMA), and a delta$+$xz numeric baseline; zstd was unavailable on the test
machine and xz/LZMA is the strong sanity check. The ablation swaps only the
residual law under a common AR predictor and coder.

\begin{table}[h]
\centering\small
\begin{tabular}{lrr|lrr}
\toprule
\multicolumn{3}{c|}{ECG (216\,kB int16)} & \multicolumn{3}{c}{S\&P (480\,kB int32)}\\
method & bytes & vs xz & method & bytes & vs xz\\
\midrule
gzip$-9$            & $118{,}619$ & $-37\%$ & gzip$-9$            & $323{,}128$ & $-23\%$\\
xz$-9$              & $86{,}772$  & ---     & xz$-9$              & $263{,}424$ & ---\\
bzip2$-9$           & $73{,}690$  & $+15\%$ & bzip2$-9$           & $254{,}343$ & $+3\%$\\
delta$+$xz          & $68{,}696$  & $+21\%$ & delta$+$xz          & $269{,}996$ & $-2\%$\\
AR$+$Gaussian       & $61{,}505$  & $+29\%$ & AR$+$Gaussian       & $244{,}229$ & $+7\%$\\
AR$+$Student-$t$    & $60{,}797$  & $+30\%$ & AR$+$Student-$t$    & $242{,}845$ & $+8\%$\\
AR$+$Stud.$+$EWMA   & $60{,}062$  & $+31\%$ & AR$+$Stud.$+$EWMA   & $242{,}216$ & $+8\%$\\
\textbf{AR$+$Mittag}& $\mathbf{59{,}686}$ & $\mathbf{+31\%}$ & \textbf{AR$+$Mittag}& $\mathbf{242{,}097}$ & $\mathbf{+8\%}$\\
AR$+$IBM mixture    & $60{,}983$  & $+30\%$ & AR$+$IBM mixture    & $314{,}849$ & $\mathbf{-20\%}$\\
\bottomrule
\end{tabular}
\caption{Actual compressed bytes, exact round-trip, all overhead included. The
codec beats every general-purpose baseline; the IBM mixture, with its fixed
light tail, \emph{fails} on power-law returns ($-20\%$, worse than gzip).}
\label{tab:stage3}
\end{table}

Three honest readings of \cref{tab:stage3}. (i) The codec is a real win: it beats
gzip, bzip2, xz, and delta$+$xz on both streams---$+31\%$ over xz on ECG, $+8\%$
on S\&P. (ii) The \emph{driver is AR$+$arithmetic coding}, not the exotic tail:
plain AR$+$Gaussian already reaches $+29\%$ / $+7\%$, and the entire heavy-tail
$+$ clock ladder adds only $\approx3\%$ (ECG) and $\approx0.9\%$ (S\&P) on top,
with the random-clock (Mittag-Leffler) model best by a hair. (iii) The
\emph{IBM mixture is not robust}: its fixed $\beta=4/3$ tail starves the frequent
large residuals of power-law returns and the code length explodes ($-20\%$); it
is only competitive where the empirical tail matches it (ECG).

\subsection{Audio: against FLAC}
\label{sec:flac}

We repeated the byte-level test on real mono $16$-bit PCM (an $11$-second
$16$\,kHz speech clip). With no \texttt{flac} binary available we implemented
FLAC's core---linear prediction plus Rice coding---as the baseline:
\textsf{FLAC-fixed} (best fixed predictor per block) and \textsf{FLAC-LPC}
(an order-$32$ linear predictor, the same family as our codec), each with
per-partition optimal Rice parameters.

\begin{table}[h]
\centering\small
\begin{tabular}{lrr}
\toprule
method & bytes & vs FLAC-LPC\\
\midrule
gzip / bzip2 / xz $-9$ & $279$--$310$\,k & $-33$ to $-48\%$\\
FLAC-fixed$+$Rice      & $229{,}944$ & $-9.8\%$\\
FLAC-LPC$32+$Rice (boss)& $209{,}340$ & ---\\
AR$+$Gaussian          & $221{,}482$ & $-5.8\%$\\
AR$+$Student-$t$       & $219{,}930$ & $-5.1\%$\\
AR$+$IBM mixture       & $220{,}270$ & $-5.2\%$\\
AR$+$Student-$t+$EWMA  & $208{,}645$ & $+0.3\%$\\
\textbf{AR$+$Mittag}   & $\mathbf{208{,}194}$ & $\mathbf{+0.5\%}$\\
\bottomrule
\end{tabular}
\caption{Lossless coding of $16$-bit speech (exact round-trip). We nudge past a
faithful FLAC-class baseline by $+0.3$ to $+0.5\%$---entirely via the volatility
model. Every \emph{static} tail (Gaussian, Student-$t$, IBM mixture) \emph{loses}
to FLAC's Rice coder by $4$--$6\%$.}
\label{tab:flac}
\end{table}

The audio result is the sharpest test of the thesis, and it is sobering. On
continuous speech we beat the general-purpose zips by $25$--$35\%$ and nudge past
FLAC-LPC by $+0.5\%$---but \emph{only} through the EWMA / Mittag-Leffler
volatility model. Every static heavy-tail code, including the IBM mixture,
\emph{loses} to FLAC's Rice coder: a tuned likelihood does not beat Rice coding
of well-predicted residuals. A real tuned \texttt{flac} (optimal LPC order and
partitioned-Rice search) would likely reclaim that sub-$1\%$, so against genuine
FLAC this is a near-tie on continuous audio, with the volatility clock the only
lever. On a quieter $8$\,kHz clip with long near-silent runs, general-purpose
bzip2/xz beat both FLAC and our codec by $40\%$---linear-prediction codecs are
simply the wrong tool when the signal is LZ-compressible.

\section{Stage 4: a learned distribution head (GPU)}
\label{sec:stage4}

A small causal GRU emits per-step distribution parameters, trained to minimize
\eqref{eq:codelen}, with a swappable Gaussian or Student-$t$ head.

\begin{table}[h]
\centering\small
\begin{tabular}{lcc}
\toprule
model & S\&P ($\bps$) & ECG ($\bps$)\\
\midrule
neural Student-$t$       & $\mathbf{1.694}$ & $\mathbf{-0.748}$\\
neural Gaussian          & $1.863$ & $-0.548$\\
classical EWMA$+$Student-$t$ & $1.789$ & $0.961$\\
\midrule
head gain (t vs.\ Gaussian) & $+0.169$ & $+0.200$\\
neural vs.\ classical clock  & $+0.095$ & $+1.71$\\
\bottomrule
\end{tabular}
\caption{Differential held-out $\bps$ (model comparison; lower is better).
The non-Gaussian head and learned volatility give the largest gains.}
\label{tab:stage4}
\end{table}

This is where the non-Gaussian thesis pays the most: a heavy-tailed head is
worth $+0.17$ to $+0.20~\bps$ over a Gaussian head \emph{at equal model size},
\emph{larger} than the static $\approx0.1$, because once a strong predictor
whitens the volatility the residual is still fat-tailed and the head captures
exactly that. The learned model also beats the hand-crafted volatility clock by
$+0.10$ (equities) to $+1.71~\bps$ (ECG). The caveat is amortization: the model
($\sim0.1$\,MB) only pays off as a pretrained, shared codec, not per small file.

\section{Repeatability of the fractional-clock advantage}
\label{sec:repeat}

The byte-level ablations showed the random-clock (Mittag-Leffler) model edging
AR$+$Student-$t+$EWMA by only $0.05$--$0.6\%$---margins small enough to be noise.
We settle this with paired tests, treating each stream as many independent
units, under two protocols. (i) \emph{Byte-level chunk sweep}: cut a stream into
independent fixed-size chunks, compress each one fully and independently
(per-chunk fit and header) with both models, and pair the resulting file sizes.
(ii) \emph{Header-free per-block}: fit once over the whole recording and compare
the two models' actual per-symbol arithmetic code length ($-\log_2$ of the coded
frequency) aggregated into $256$-sample blocks---removing per-chunk header and
fit noise to isolate the model. Both are run on \emph{two independent} ECG
recordings: the SciPy electrocardiogram and MIT-BIH record~100 (a different
subject and lead, $650$k samples).

\begin{table}[h]
\centering\small
\begin{tabular}{lrrl}
\toprule
stream (header-free, per block) & blocks & mean impr. & win rate / $p$\\
\midrule
ECG (SciPy)        & $421$  & $+0.64\%$ & $96.2\%$ / $10^{-68}$\\
ECG (MIT-BIH 100)  & $2539$ & $+0.80\%$ & $99.8\%$ / ${\approx}0$\\
S\&P returns       & $75$   & $+0.009\%$ & significant but trivial\\
audio (speech)     & $21$   & $+0.05\%$ & $57\%$ / $0.18$ (noise)\\
\bottomrule
\end{tabular}
\caption{Paired per-block code-length comparison, Mittag-Leffler vs.\
Student-$t+$EWMA. The advantage is robust and reproduces across two independent
ECG recordings; it is statistically real but practically zero on returns and
indistinguishable from noise on audio.}
\label{tab:repeat}
\end{table}

At the model level the result is unambiguous on ECG: on both recordings the
fractional kernel wins $96$--$99.8\%$ of all blocks with overwhelming
significance, because it matches the long-memory volatility that EWMA's
exponential decay cannot (\cref{sec:stage2}). At the byte level the same
advantage appears but is \emph{conditional on block size}: with per-chunk
headers it \emph{loses} on $1$\,kB chunks (the dynamic model's header overhead
and its $128$-lag warmup are not amortized), turns positive by $\sim2$--$4$\,kB,
and reaches $+0.35$ to $+0.52\%$ at $8$\,kB---with the same pattern on both
recordings.

\begin{finding}
The Mittag-Leffler / fractional-clock advantage over EWMA is real and
reproducible across independent ECG recordings ($+0.6$--$0.8\%$ at the model
level, $96$--$99.8\%$ of blocks), but small, restricted to
long-memory-volatility data, and---at the byte level---realized only once the
coding block is large enough to amortize the model's overhead. It is negligible
on returns and noise on audio.
\end{finding}

\section{Verdict}
\label{sec:verdict}

\begin{itemize}
\item \textbf{A real codec, byte-for-byte.} With exact round-trip and all
  overhead included, the AR$+$arithmetic codec beats every general-purpose
  baseline---gzip, bzip2, and xz---by $+8\%$ (returns) to $+31\%$ (ECG) over
  xz, and it \emph{ties} a faithful FLAC-class baseline on continuous speech
  ($+0.5\%$, within real-FLAC tuning margin). The structure compresses.
\item \textbf{But the driver is prediction $+$ coding, not the exotic tail.}
  Plain AR$+$Gaussian already captures almost the entire margin; the whole
  heavy-tail $+$ clock ladder adds only $\approx1$--$3\%$ on numeric data, and on
  audio every \emph{static} heavy-tail code (including the IBM mixture)
  \emph{loses} to FLAC's Rice coder. The one consistent lever above a Gaussian
  baseline is the volatility model, not the likelihood shape---and even its best
  (fractional) form beats EWMA repeatably only on long-memory ECG, by under
  $1\%$ (\cref{sec:repeat}).
\item \textbf{The specific operator predictions do not transfer, and can hurt.}
  Real tails are power-law (Student-$t$), not the biharmonic $\beta=4/3$; the
  fractional volatility kernel confirms long memory but beats GARCH only
  marginally; and the \emph{IBM mixture fails outright} on power-law returns
  ($-20\%$ vs xz, worse than gzip). The raw order-$p$ pseudo-process kernels are
  signed---not codes at all; only their positive subordinated (IBM) face
  compresses, and only where the empirical tail happens to match it.
\item \textbf{The framework is the contribution.} What survives contact with real
  data is the architecture the papers motivate: an adaptive heavy-tail head
  (selected per block), a long-memory volatility clock, and the view---made
  precise by subordination---that tail and volatility are one mechanism, to be
  modeled jointly rather than bolted together.
\end{itemize}

A deployable artifact consistent with these results is a per-domain,
\emph{pretrained neural compressor with a heavy-tailed (Student-$t$) head and a
long-memory volatility state}; it captures essentially all the measured gains.
The iterated-Brownian and order-$p$ theory earns its place as the conceptual
scaffolding for that design, even where its closed-form exponents are not what
the data wants.

\section*{Reproducibility}
All experiments are pure Python (NumPy/SciPy; PyTorch with CUDA for
Stage~4) and reproducible from the scripts \texttt{stage0\_synthetic.py},
\texttt{memdiag.py}, \texttt{memvol.py}, \texttt{disentangle.py},
\texttt{stage1\_realdata.py}, \texttt{stage2\_mittag.py}, \texttt{codec.py},
\texttt{pseudo\_codes.py} (the genuine FFT-built order-$p$ kernels and the IBM
mixture code), \texttt{bench.py} and \texttt{bench\_wav.py} (the byte-level
codec with exact round-trip and the gzip/bzip2/xz and FLAC-class baselines),
\texttt{repeat\_test.py}, \texttt{ecg\_repeat.py} and \texttt{mit\_test.py} (the
paired repeatability and cross-recording tests), and \texttt{stage4\_neural.py}. The ECG is the SciPy electrocardiogram; the equity
series is pooled S\&P~500 daily log-returns. The arithmetic coder is
round-trip verified lossless.

\begin{thebibliography}{9}
\bibitem{AllouZheng} H.~Allouba and W.~Zheng,
  \emph{Brownian-time processes: the PDE connection and the half-derivative
  generator}, Ann.\ Probab.\ \textbf{29} (2001), 1780--1795.
\bibitem{chaos} R.~Fontes,
  \emph{Source-Resolved Chaos Expansions for Iterated Brownian Motion}, 2026.
\bibitem{ito} R.~Fontes,
  \emph{Iterated Brownian Motion: An It\^o Formula in the Energy-First
  Framework}, 2026.
\bibitem{orderp} R.~Fontes,
  \emph{A Level-$p$ Clark--Ocone Calculus for Higher-Order Heat
  Pseudo-Processes}, 2026.
\bibitem{variance} R.~Fontes,
  \emph{A Numerical Variance Budget for Iterated Brownian Motion}, 2026.
\bibitem{mem} R.~Fontes,
  \emph{The Fractional Memory Process: Well-Posedness, Spectral Decomposition,
  and Discrete Representation}, 2026.
\end{thebibliography}

\end{document}
